% Table presentation generated by scripts/editorial_revision.py.
\begin{table}[!htbp]
\centering
\begin{threeparttable}
\caption{Mean forecast scores and skill relative to seasonal climatology across five origins.}
\label{tab:comment2_external_comparison}
\fontsize{8.5}{10.5}\selectfont
\begin{tabular}{lrrrr}
\hline
Method & 28-day CRPS & 28-day CRPSS & 8-day CRPS & 8-day CRPSS \\
\hline
Seasonal climatology & 0.64218 & 0.00000 & 0.72931 & 0.00000 \\
Persistence & 0.91061 & -0.41801 & 0.99854 & -0.36915 \\
GloFAS & 0.61761 & 0.03826 & 0.77502 & -0.06267 \\
NWS & --- & --- & 0.66478 & 0.08848 \\
GloFAS EMOS & 0.45734 & 0.28782 & 0.44450 & 0.39052 \\
AL-M-T1 & 0.23230 & 0.63826 & 0.30944 & 0.57570 \\
exAL-M-T1 & 0.20164 & 0.68601 & 0.29926 & 0.58967 \\
\hline
\end{tabular}
\begin{tablenotes}
\fontsize{8}{10}\selectfont
\item \textit{Note:} Daily USGS discharge supplies verification. Scores use $\log(1+x)$ discharge, with $x$ in m$^3$/s. CRPS denotes the continuous ranked probability score; lower values are better. CRPSS is one minus the ratio of mean CRPS to the seasonal-climatology mean on the same cases; higher values are better. Means include 140 cases over leads 1--28 and 40 over leads 1--8. Persistence is a point forecast; GloFAS EMOS is a single-product ensemble model output statistics correction. NWS is available through lead eight. N, AL and exAL denote normal, asymmetric Laplace and extended asymmetric Laplace likelihoods. U and M denote one and multiple observation sources; T1 retains the transfer contribution at forecast leads and T0 omits it. CRPS uses quantile-grid integration as defined in Subsection~\ref{subsec:modelselection}.
\end{tablenotes}
\end{threeparttable}
\end{table}
